\section{Reconstructing training state across a checkpoint restart}
\label{sec:appendix-restart}

The unfrozen FreeMatch arm continues from its frozen-comparison head. Those checkpoints store
four keys, \restartKeys{}, so the model survives and the training state around it does not.
Continuing from such a checkpoint restarts three things: the self-adaptive threshold state,
which re-initialises to $1/C$; the consistency weight's warm-up, which holds the consistency term at
zero for the first \restartMuZeroSteps{} steps of the \restartScheduleStepsNs{}-step schedule;
and AdamW's moment estimates. The fixes below implement the pre-registered intent that the
decision boundary and the threshold start consistent with each other. Every TPR value in this
appendix is on the original test pool, as each run scored it during training.

\paragraph{The evidence.} A \restartControlStepsNs{}-step frozen continuation of the
development-fold head, with none of this state rebuilt, ran \restartStart{} $\to$
\restartTrough{} at step \restartTroughStep{} $\to$ \restartEnd{} in TPR@$10^{-3}$, a maximum
drift of \restartMaxDrift{} (\restartMaxDriftPos{} positives against a one-positive resolution
of \restartOnePos{}), with the trough inside the window where the consistency term is held at zero.

\paragraph{The fixes.} \emph{Threshold state:} with the weights fixed, the moving average has
the pool mean as its fixed point, so we rebuilt it by averaging per-batch statistics over
\restartSatPasses{} weakly augmented passes of the \restartPool{}-image target pool
(\restartSatBatches{} batches). Driving the real update rule, momentum \restartMomentum{}, for
\restartEmaItersNs{} iterations agrees with this to \restartTauAgree{} on the global threshold and
\restartPbarAgree{} on the class marginal, against a tolerance of \restartSatTol{}.
\emph{Warm-up offset:} we resume the consistency-weight schedule at step \restartMuOffsetNs{}, so
the consistency term is at full weight from the first step. \emph{Moment warm-up:} AdamW runs
\restartAdamSteps{} steps under the corrected objective with the weights restored after every
step and the encoder frozen. A probe continues to step \restartAdamProbeEnd{}, over which the
bias-corrected second moment moved \restartAdamDebiased{} from step \restartAdamSteps{}, and
then rolls the optimiser state back to step \restartAdamSteps{}; it does not roll back the
threshold state, which the probe's forward passes keep updating, or the half-precision loss
scale. We call this reconstruction, not restoration: the original values were never stored.
The order is load the head, rebuild the threshold state, warm the moments, unfreeze, train.

\paragraph{The reconstructed control.} With the three fixes, the same frozen continuation
drifts at most \restartFixedMaxDrift{} from step 0 (\restartFixedMaxDriftPos{} positives). Over
its \restartFixedN{} later checkpoints its mean is \restartFixedMean{} (sd \restartFixedSD{}),
against \restartTailMean{} (sd \restartTailSD{}) over the last \restartTailN{} checkpoints of
the frozen run itself; we ran no test of their equality.

\paragraph{The step-0 gate.} Every fold is gated before training (Table~\ref{tab:restartgate2}).
We first required exact reproduction of TPR@$10^{-3}$. That held on the development fold and on
fold 1 and refused fold 2, whose step-0 value differed by one positive: scoring through the
encoder inside the training loop drifts slightly from the cached scoring path. We recalibrated
the rule before rerunning fold 2, to TPR within one positive together with an AUROC check at
$10^{-3}$, and folds 2--4 passed.

% GENERATED FILE -- do not edit by hand.
% Rebuild:  python make_all.py (generator: make_appendix_relabel_v2.py)
% Relabelled copy of tables/A5_gate.tex; tabular number tokens asserted identical, caption/notes token changes declared.
\begin{table}[t]
\centering
\caption{\textbf{The per-fold step-0 gate, as each run recorded it} (unfrozen FreeMatch). The gate re-scores the loaded head before any training and aborts the fold on a mismatch. Folds 0 and 1 ran under the original rule, exact reproduction of TPR@$10^{-3}$; fold 2 differed by exactly one positive, which that rule refuses, and folds 2--4 ran under the recalibrated rule. TPR differences are in positives; the AUROC half was added with the recalibration.}
\label{tab:restartgate2}
\footnotesize
\setlength{\tabcolsep}{4pt}
\begin{tabular}{llrr}
\toprule
Fold & Rule & $|\Delta\mathrm{TPR}|$ (positives) & $|\Delta\mathrm{AUROC}|$ \\
\midrule
0 (development) & TPR exact & 0 & --- \\
1 (confirmatory) & TPR exact & 0 & --- \\
2 (confirmatory) & TPR $\le$ 1 positive; AUROC $\le 10^{-3}$ & 1 & $3.7 \times 10^{-6}$ \\
3 (confirmatory) & TPR $\le$ 1 positive; AUROC $\le 10^{-3}$ & 0 & $9.5 \times 10^{-6}$ \\
4 (confirmatory) & TPR $\le$ 1 positive; AUROC $\le 10^{-3}$ & 0 & $1.4 \times 10^{-5}$ \\
\bottomrule
\end{tabular}
\end{table}

